
% ==================================================
% 1. CLASSE DEL DOCUMENT
% ==================================================
\documentclass[a4paper,12pt,oneside]{article}

% ==================================================
% 2. PAQUETS ESSENCIALS (IDIOMA I CODIFICACIÓ)
% ==================================================
\usepackage[utf8]{inputenc}
\usepackage[T1]{fontenc}
\usepackage[catalan]{babel}
\usepackage{lmodern}

% ==================================================
% 3. FORMAT DEL TEXT I DELS PARÀGRAFS
% ==================================================
\setlength{\parindent}{0.5cm}
\usepackage{indentfirst}

\usepackage{setspace}
\singlespacing

% ==================================================
% 4. MATEMÀTIQUES
% ==================================================
\usepackage{amsmath,amssymb,amsthm}
\usepackage{mathtools}
\usepackage{bm}
\usepackage{siunitx}

% ==================================================
% 5. DISSENY I FORMAT DE LA PÀGINA
% ==================================================
\usepackage{geometry}
\geometry{
    a4paper,
    margin=2.5cm,
    top=2.5cm,
    bottom=2.5cm
}

\raggedbottom

% ==================================================
% 6. TÍTOLS I SECCIONS
% ==================================================
\usepackage{titlesec}

% Format de les seccions principals
\titleformat{\section}
  {\normalfont\Large\bfseries}
  {\thesection.}
  {0.5em}
  {}

% Format de les subseccions
\titleformat{\subsection}
  {\normalfont\large\bfseries}
  {\thesubsection.}
  {0.5em}
  {}

% Format de les subsubseccions
\titleformat{\subsubsection}
  {\normalfont\normalsize\bfseries}
  {\thesubsubsection.}
  {0.5em}
  {}

% Espaiats dels títols
\titlespacing*{\section}{0pt}{18pt}{10pt}
\titlespacing*{\subsection}{0pt}{12pt}{6pt}
\titlespacing*{\subsubsection}{0pt}{8pt}{4pt}

% ==================================================
% 7. CAPÇALERES I PEUS DE PÀGINA
% ==================================================

% ==================================================
% CAPÇALERES I PEUS DE PÀGINA
% ==================================================

\usepackage{fancyhdr}

\pagestyle{fancy}
\fancyhf{}

% Mostrar el número i el nom de la secció
\renewcommand{\sectionmark}[1]{%
    \markboth{\thesection.\ #1}{}
}

% Capçalera: nom de la secció a l'esquerra
\fancyhead[L]{\nouppercase{\leftmark}}

% Número de pàgina centrat al peu
\fancyfoot[C]{\thepage}

% Línies de capçalera i peu
\renewcommand{\headrulewidth}{0.4pt}
\renewcommand{\footrulewidth}{0pt}

% Configuració de les pàgines plain
\fancypagestyle{plain}{
    \fancyhf{}
    \fancyhead[L]{\nouppercase{\leftmark}}
    \fancyfoot[C]{\thepage}
    \renewcommand{\headrulewidth}{0.4pt}
    \renewcommand{\footrulewidth}{0pt}
}

% ==================================================
% 8. GRÀFICS, TAULES I COLORS
% ==================================================
\usepackage{graphicx}
\usepackage[table]{xcolor}
\usepackage{array}
\usepackage{tabularx}
\usepackage{multirow}
\usepackage{booktabs}
\usepackage{makecell}
\usepackage{float}
\usepackage[font=small]{caption}
\usepackage{subcaption}

% Caixes de text
\usepackage{tcolorbox}
\tcbuselibrary{most}

% Colors
\definecolor{grisC}{RGB}{211,211,211}
\definecolor{grisF}{gray}{0.3}

% ==================================================
% 9. ENLLAÇOS I REFERÈNCIES
% ==================================================
\usepackage[hidelinks]{hyperref}
\usepackage{bookmark}

\usepackage[autostyle]{csquotes}
\DeclareQuoteStyle{catalan}
  {\guillemotleft}{\guillemotright}
  {\quotesinglbase}{\guillemotright}

\usepackage[
    backend=biber,
    style=numeric,
    sorting=none
]{biblatex}

% ==================================================
% 10. ALTRES PAQUETS ÚTILS
% ==================================================
\usepackage{varwidth}
\usepackage{enumitem}
\usepackage[title,titletoc]{appendix}

% ==================================================
% 11. CONFIGURACIÓ D'ENTORNS DE TEOREMES
% ==================================================
% \theoremstyle{definition}
% \newtheorem{definicio}{Definició}[section]
%
% \theoremstyle{plain}
% \newtheorem{teorema}{Teorema}[section]
% \newtheorem{propietat}{Propietat}[section]
%
% \theoremstyle{remark}
% \newtheorem{exemple}{Exemple}[section]
% \newtheorem{nota}{Nota}[section]

% ==================================================
% 12. INICI DEL DOCUMENT
% ==================================================
\begin{document}

\begin{titlepage}
    \thispagestyle{empty}
    \centering

    % Universitat
    \vspace*{1.5cm}

    {\Large\scshape Universitat Autònoma de Barcelona\par}
    \vspace{0.3cm}
    {\large Grau en Física\par}

    % Títol principal
    \vspace{2.5cm}

    {\color{gray}\rule{1\textwidth}{1.2pt}\par}

    \vspace{0.8cm}

    {\large\bfseries PRÀCTICA 1\par}

    \vspace{0.5cm}

    {\Huge\bfseries Modelització d'una ablació cardíaca IVS\par}

    \vspace{0.6cm}

    {\large INFORME DE PRÀCTIQUES\par}

    \vspace{0.8cm}

    {\color{gray}\rule{1\textwidth}{1.2pt}\par}

    % Dades dels grups
    \vspace{2cm}

    {\large\bfseries Grup 111\par}

    \vspace{0.8cm}

    {\large Nom i NIU\par}
    \vspace{0.2cm}
    {\large Aleix Fosalba Pujols 1709441\par}
    \vspace{0.2cm}
    {\large Victor Shenyu Lechuga Izquierdo 1753089\par}
    \vspace{0.2cm}
    

    % Informació inferior
    \vfill

    {\large\bfseries Mètodes Numèrics II\par}

    \vspace{0.3cm}

    {\large Curs acadèmic 2026--2027\par}

    \vspace{1cm}

\end{titlepage}

\newpage

\tableofcontents

\newpage

\section{Introducció}

L'ablació cardíaca és un procediment mèdic utilitzat per tractar arrítmies cardíaques mitjançant l'aplicació de calor sobre el 
teixit responsable de les alteracions elèctriques. En aquesta pràctica s'estudia l'ablació cardíaca interventricular (IVS), 
basada en l'aplicació d'una diferència de potencial entre dos elèctrodes de polaritat oposada. Aquest procés genera calor per 
efecte Joule, provocant un augment de la temperatura del teixit i permetent destruir les cèl·lules afectades.

L'objectiu principal és determinar el temps òptim d'aplicació d'una tensió alterna de 40 V $\text{(pi a pic)}$ sobre un teixit 
cardíac de 2 cm de gruix, amb una regió central de 0.5 cm corresponent al teixit malalt. El tractament ha de garantir que aquesta 
regió assoleixi temperatures superiors als 50 $^\circ$C, evitant que el teixit sa superi aquest llindar i que cap zona arribi als 80 $^\circ$C, 
temperatura a partir de la qual es poden produir fenòmens d'evaporació i altres complicacions mèdiques.

Per estudiar l'evolució temporal i espacial de la temperatura, es modelitza la transferència de calor mitjançant l'equació de 
Fourier en una dimensió, Eq. (\ref{fourier}), incorporant la generació de calor per efecte Joule. La resolució d'aquesta equació permet 
simular el procés i analitzar-ne l'eficàcia sense haver de realitzar la intervenció. 

Finalment, es resol el model mitjançant tres mètodes numèrics: Euler explícit, Euler implícit i Crank-Nicolson. Es comparen els 
resultats obtinguts per als diferents mallats, estudiant-ne l'error numèric, l'estabilitat i la precisió, amb la finalitat de 
determinar les condicions més adequades per dur a terme el tractament.

\section{Mètodes numèrics}

\subsection{Normalització}

La llei de Fourier per la temperatura és:

\begin{equation}
    c_v \rho \dfrac{\partial T}{\partial x} = \nabla (\kappa \nabla T) + P_{\text{ext}},
    \label{fourier}
\end{equation}

on $c_v$ és el calor específic, $\rho$ és la densitat, $\kappa$ és la conductivitat tèrmica i $P_{\text{ext}}$ representa totes 
les fonts de calor externes. Per aquest problema considerarem que l'única font de calor externa és la calor produïda pels 
elèctrodes que generen calor per efecte Joule.

A continuació definim les variables normalitzades:

\begin{equation}
    \widetilde{T}=\frac{T}{T_0};\quad \widetilde{x}=\frac{x}{x_0};\quad \widetilde{t}=t/t_0
    \label{norm}
\end{equation}
d'on trobem l'equació:
\begin{equation*}
    \frac{\partial \widetilde{T}}{\partial \widetilde{t}} = \frac{\partial^2 \widetilde{T}}{\partial \widetilde{x}^2}  
    \frac{t_0\kappa}{x_0^2c_v\rho}+\frac{\sigma (\Delta V)^2t_0}{T_0d^2c_v\rho}
\end{equation*}
I en particular:
 \begin{equation}
     x_0=d \quad t_0=\frac{d^2}{\alpha} \quad T_0=\frac{\sigma (\Delta V)^2}{\kappa}
\end{equation}
Ara ja podem normalitzar l'eq (\ref{fourier}). La llei de Fourier normalitzada és:

\begin{equation}
    \frac{\partial \widetilde{T}}{\partial \widetilde{t}}=\frac{\partial^2 \widetilde{T}}{\partial \widetilde{x}^2} + 1
\end{equation}

Els valors numèrics emprats en aquesta pràctica són: 
\begin{equation*}
    c_v = 3698 \ \frac{\text{J}}{\text{kg \ K}}, \quad \rho = 1081 \ \frac{\text{kg}}{\text{m} ^3}, \quad \kappa = 0.56 \ \frac{\text{W}}{\text{m \ K}}, \quad \sigma = 0.472 \ \frac{\text{S}}{\text{m}}
\end{equation*}
\begin{equation*}
    x_0 = 0.02 \ \text{m}, \quad t_0 approx 2855.384 \ \text{s}, \quad T_0 \approx 674.286 \ \text{K}
\end{equation*}

\subsection{Discretització del sistema}

La discretització del sistema difereix en els mètodes aplicats en aquesta pràctica.

En el cas del mètode d'Euler explícit, considerem un mallat espacial amb un nombre de punts $N$ = 101 i d'amplada $\Delta \widetilde{x} = 0.01$ (normalitzat). 
Pel mallat temporal, considerem aquests tres (també normalitzats):
\begin{equation*}
    \Delta \widetilde{t} = 0.51 (\Delta \widetilde{x})^2 = 5.1 \cdot 10^{-5}, \quad \Delta \widetilde{t} = 0.49 (\Delta \widetilde{x})^2 = 4.9 \cdot 10^{-5}, \quad \Delta \widetilde{t} = 0.25 (\Delta \widetilde{x})^2 = 2.5 \cdot 10^{-5}
\end{equation*}

Amb un nombre de punts temporals iguals a 490, 510 i 1000, respectivament.

Per l'altra banda, en el cas dels mètodes d'Euler implícit i Crank-Nicolson, considerem el mateix mallat espacial que a l'Euler 
implícit. En el cas de la discretització temporal, estudiem els següents dos mallats:
\begin{equation*}
    \Delta \widetilde{t} = (\Delta \widetilde{x})^2, \quad \Delta \widetilde{t} = 0.5 (\Delta \widetilde{x})^2
\end{equation*}

Amb un nombre de punts temporals iguals a 250 i 500, respectivament.

\subsection{Condicions inicials i de contorn}

El problema es defineix amb una equació diferencial en derivades parcials. Per solucionar aquest problema cal conèixer les condicions 
inicials i les condicions de contorn a les fronteres. Volem modelitzar una operació cardíaca de manera que considerem la temperatura 
inicial com la temperatura mitjana del cos humà, uns $T_c = 36.5 ^\circ$C. A més, com que el teixit de l'estudi es troba a dins 
del cor i envoltat per grans quantitats de sang en circulació contínua, podem considerar que la temperatura a les fronteres és 
la mateixa que la del cos, és a dir, $T_c$.


\subsection{Mètode d'Euler explícit}

\subsection{Mètode de Crank Nicolson}


Per resoldre numèricament l’equació del calor adimensionalitzada s’ha implementat el mètode de Crank-Nicolson en Python. Aquest mètode consisteix a aproximar la derivada temporal mitjançant una diferència finita i la derivada espacial de segon ordre mitjançant la mitjana de les aproximacions als instants temporals actual i posterior.

S’ha discretitzat el domini espacial normalitzat entre 0 i 1 utilitzant 101 punts, amb una separació espacial de $(\Delta x=0.01)$. S’han imposat les condicions inicials i de contorn corresponents a una temperatura constant de $36.5\,^\circ\mathrm{C} $.

La discretització de l’equació diferencial dona lloc a un sistema d’equacions lineals que s’ha expressat matricialment. Per resoldre’l a cada pas temporal s’ha utilitzat la factorització LU, implementada mitjançant la llibreria SciPy.

S’han estudiat dos passos temporals diferents, $(\Delta t=(\Delta x)^2)$ i $(\Delta t=0.5(\Delta x)^2)$, fins a arribar al temps adimensional (t=0.025).

Finalment, s’han comparat les solucions numèriques amb la solució analítica obtinguda mitjançant una sèrie de Fourier, calculant l’error absolut, l’error màxim i l’error quadràtic mitjà. Aquesta comparació permet avaluar la precisió del mètode i estudiar la influència del pas temporal sobre els resultats.

\subsection{Anàlisi de resultats}
\textbf{QUEREIS PONER ESTA PARTE DE ANALISI DE RESULTADOS COMO UNA SUBSECTION A PARTE? O EN CADA METODO HACEMOS UN ANALISI DE RESULTADOS EN VEZ DE HACER DESPUES UN APARTADO CON TODOS? (NOSE SI ME HE EXPLICADO BIEN)}
Les solucions numèriques obtingudes amb els dos passos temporals estudiats presenten una gran coincidència amb la solució analítica, tal com s’observa a la primera gràfica.

Respecte a l’error numèric, s’obtenen errors màxims de l’ordre de $(10^{-6})$ en ambdós casos. S’observa que la reducció del pas temporal no comporta una disminució significativa de l’error. Això indica que, per a aquesta discretització, la contribució de l’error espacial és predominant.

Finalment, la distribució de temperatura en graus Celsius mostra que la temperatura màxima al centre del teixit és aproximadament de $53.17\,^\circ\mathrm{C}$, mentre que als extrems es manté a $36.5\,^\circ\mathrm{C}$. Per tant, la temperatura supera els $50\,^\circ\mathrm{C}$ a la regió central, però no arriba als $80\,^\circ\mathrm{C}$ en cap punt del domini.

Aquests resultats indiquen que, en el temps estudiat, s’assoleix el mínim de temperatura necessària per provocar la mort cel·lular sense arribar al límit d’evaporació.





\end{document}